\documentclass{article}
\usepackage{geometry}
\geometry{
	a4paper,
	left=2.5cm,      % 左边距 2cm
	right=2.5cm,     % 右边距 2cm
	top=2.5cm,       % 上边距 2cm
	bottom=2.5cm,    % 下边距 2cm
	headheight=13.6pt
}
\usepackage{amsmath}
\usepackage{booktabs} % 用于三线表
\usepackage{caption}
\usepackage{setspace} % 用于设置行间距
\usepackage{array}
\usepackage{graphicx} % 用于插入图片
\setlength{\headheight}{13.6pt}
\usepackage{enumitem} % 用于自定义列表格式
\usepackage{svg}

\begin{document}
	\subsection{Using Spearman Correlation Coefficient to Determine Category Correlations}
	
	\subsubsection{Spearman Correlation Coefficient Algorithm}
	
	The Spearman \textbf{correlation coefficient} is widely used to measure the strength and direction of \textbf{monotonic relationships} between two variables. Its value ranges from \([-1,1]\). A positive coefficient indicates a positive monotonic correlation between the two variables, while a negative coefficient indicates a negative monotonic correlation. The closer the absolute value is to 1, the stronger the monotonic correlation between the variables; the closer it is to 0, the weaker the monotonic correlation. The calculation steps are as follows:
	
	\begin{enumerate}[label=\arabic*., leftmargin=*]
		\item \textbf{Convert data to ranks}
		
		Rank all values of variable \(X\) in ascending order, with rank 1 being the minimum value and rank \(n\) being the maximum value. If there are tied values, assign them the average of their ranks. Perform the same ranking for variable \(Y\).
		
		\item \textbf{Calculate the Spearman correlation coefficient between the two rank sequences}
		
		Substitute the ranked values into the formula:
		\begin{equation}
			\rho = \frac{\sum_{i=1}^{n} \left(R_{x_i} - \overline{R_x}\right) \left(R_{y_i} - \overline{R_y}\right)}
			{\sqrt{\sum_{i=1}^{n} \left(R_{x_i} - \overline{R_x}\right)^2 \sum_{i=1}^{n} \left(R_{y_i} - \overline{R_y}\right)^2}}
			\label{eq:spearman_correlation}
		\end{equation}
		where \(\rho\) is the Spearman correlation coefficient, $\bar{R}_x$ is the mean of $R_x$ and $\bar{R}_y$ is the mean of $R_y$.
	\end{enumerate}
	
	\subsubsection{Solution Results}
	
	The Spearman correlation coefficient was calculated using MATLAB, and a heatmap was generated:
	
	\begin{figure}[htbp]
		\centering
		\includegraphics[width=0.8\textwidth]{figures/pearson_picture.eps}
		\caption{Spearman Correlation Coefficient Heatmap for Category Sales}
		\label{fig:spearman_heatmap}
	\end{figure}
	
	Based on the above table, no strongly correlated categories were identified. The relatively strongly correlated categories are listed in the table below:
	
	\begin{table}[htbp]
		\centering
		\caption{Relatively Strongly Correlated Category Combinations (\(|\rho|>0.5\))}
		\label{tab:spearman_strong_corr}
		\renewcommand{\arraystretch}{1.5}
		\renewcommand{\heavyrulewidth}{1.2pt}
		\begin{tabular}{ccc}
			\toprule
			Category 1 & Category 2 & Correlation Coefficient \\
			\midrule
			Leafy Vegetables & Cauliflower & 0.638288645480482 \\
			Leafy Vegetables & Aquatic Stem Vegetables & 0.568835111289716 \\
			Leafy Vegetables & Peppers & 0.668577433074466 \\
			Leafy Vegetables & Edible Mushrooms & 0.639654735509754 \\
			Cauliflower & Aquatic Stem Vegetables & 0.549840658198361 \\
			Cauliflower & Peppers & 0.561693690138146 \\
			\bottomrule
		\end{tabular}
		\renewcommand{\heavyrulewidth}{0.8pt}
	\end{table}
	
	The above relatively strongly correlated category combinations are also presented in the file Correlated Categories.xlsx.
	
	\subsubsection{Spearman Correlation Coefficient Test}
	
	This study conducted a Spearman correlation coefficient test on the above results, following the steps below:
	
	\paragraph{Step 1. Propose Hypotheses}
	
	Null hypothesis \(H_0\): The true correlation coefficient \(\rho = 0\), meaning there is no monotonic correlation between the two variables in the population.
	
	Alternative hypothesis \(H_1\): The true correlation coefficient \(\rho \neq 0\), meaning there is a monotonic correlation between the two variables in the population.
	
	The confidence level is set to 95\%.
	
	\paragraph{Step 2. Calculate the \(P\) Value}
	
	The p-value is derived from a \textbf{t-statistic} transformed from \(r\). The formula for the t-statistic is as follows:
	
	\begin{equation}
		t = \frac{r - \rho_0}{SE_r} = \frac{r}{\sqrt{\frac{1-r^2}{n-2}}}
		\label{eq:t_test_correlation}
	\end{equation}
	
	Where \(r\) represents the calculated sample Spearman correlation coefficient, \(\rho_0\) represents the population correlation coefficient under the null hypothesis, \(SE_r\) represents the standard error of the correlation coefficient \(r\), and \(n-2\) represents the degrees of freedom (df). Since the null hypothesis is \(\rho = 0\), the formula for the t-statistic can be simplified to:
	
	\begin{equation}
		t = r \sqrt{\frac{x-2}{1-r^2}}
		\label{eq:t_test_simplified}
	\end{equation}
	
	The calculated t-statistic follows a \textbf{t-distribution} with \(\textit{df} = n - 2\) degrees of freedom. The p-value is the sum of the areas in both tails of the t-distribution.
	
	This study used MATLAB to calculate the sum of the areas in both tails, performing the Spearman correlation coefficient test. The results are shown in the tables below:
	
	\begin{table}[htbp]
		\centering
		\caption{\(P\) Values for Spearman Correlation Coefficient Test of Relatively Strongly Correlated Categories}
		\label{tab:spearman_p_strong}
		\renewcommand{\arraystretch}{1.5}
		\renewcommand{\heavyrulewidth}{1.2pt}
		\begin{tabular}{ccc}
			\toprule
			Category 1 & Category 2 & \(P\) Value \\
			\midrule
			Leafy Vegetables & Cauliflower & 0.0000000 \\
			Leafy Vegetables & Peppers & 0.0000000 \\
			Cauliflower & Eggplant & 0.0000000 \\
			Cauliflower & Peppers & 0.0000000 \\
			Eggplant & Peppers & 0.0000000 \\
			Peppers & Edible Mushrooms & 0.0000000 \\
			\bottomrule
		\end{tabular}
		\renewcommand{\heavyrulewidth}{0.8pt}
	\end{table}
	
	\begin{table}[htbp]
		\centering
		\caption{\(P\) Values for Spearman Correlation Coefficient Test Across All Categories}
		\label{tab:spearman_p_all}
		\renewcommand{\arraystretch}{1.5}
		\renewcommand{\heavyrulewidth}{1.2pt}
		\begin{tabular}{ccccccc}
			\toprule
			& Leafy Vegetables & Cauliflower & Aquatic Stem & Eggplant & Peppers & Edible Mushrooms \\
			\midrule
			Leafy Vegetables & 0.0000000 & 0.0000000 & 0.0000000 & 0.0000000 & 0.0000000 & 0.0000000 \\
			Cauliflower & 0.0000000 & 0.0000000 & 0.0000000 & 0.0000000 & 0.0000000 & 0.0000000 \\
			Aquatic Stem & 0.0000000 & 0.0000000 & 0.0000000 & 0.0030232 & 0.0000000 & 0.0000000 \\
			Eggplant & 0.0000000 & 0.0000000 & 0.0030232 & 0.0000000 & 0.0000000 & 0.0000051 \\
			Peppers & 0.0000000 & 0.0000000 & 0.0000000 & 0.0000000 & 0.0000000 & 0.0000000 \\
			Edible Mushrooms & 0.0000000 & 0.0000000 & 0.0000000 & 0.0000051 & 0.0000000 & 0.0000000 \\
			\bottomrule
		\end{tabular}
		\renewcommand{\heavyrulewidth}{0.8pt}
	\end{table}
	
	Based on \textbf{Table 2} and \textbf{Table 3}, the Spearman correlation coefficients between the sales volumes of all categories are less than the confidence level of 0.05. Therefore, the null hypothesis is rejected, indicating that the conclusions drawn in the previous step are credible.
\end{document}